% 第3章 3.1.1：标准正交序列在单位算子与 A(x_n)=(x_n/n) 下的输出。
% (a) 斜投影示意 e_1,e_2,e_3 与 Ae_1=e_1, Ae_2=e_2/2, Ae_3=e_3/3；两两距离 ||e_n-e_m||=sqrt2。
% (b) 相邻输出距离：单位算子恒为 sqrt2≈1.4142；A 下为 sqrt(1/n^2+1/(n+1)^2)，n=1..15
%     （1.1180, 0.6009, 0.4167, ..., 0.0914，由 python 直接计算）。
\begin{tikzpicture}[font=\figurefont\footnotesize]
\begin{scope}[scale=1.55]
  \coordinate (O) at (0,0);
  \coordinate (E1) at (1.2,0);
  \coordinate (E2) at (0,1.2);
  \coordinate (E3) at (-0.62,-0.55);
  \draw[FigureGrid,line width=.4pt] (O) -- ($1.12*(E1)$);
  \draw[FigureGrid,line width=.4pt] (O) -- ($1.12*(E2)$);
  \draw[FigureGrid,line width=.4pt] (O) -- ($1.12*(E3)$);
  \draw[FigureBlue,dashed,line width=.7pt] (E1) -- node[above right=-1pt,text=FigureBlue] {$\sqrt2$} (E2);
  \draw[FigureBlue,dashed,line width=.7pt] (E2) -- node[left,text=FigureBlue] {$\sqrt2$} (E3);
  \draw[FigureBlue,dashed,line width=.7pt] (E3) -- node[below,text=FigureBlue] {$\sqrt2$} (E1);
  \draw[FigureBlue,line width=1pt,-{Stealth[length=1.8mm]}] (O) -- (E1) node[right,text=FigureInk] {$e_1$};
  \draw[FigureBlue,line width=1pt,-{Stealth[length=1.8mm]}] (O) -- (E2) node[above,text=FigureInk] {$e_2$};
  \draw[FigureBlue,line width=1pt,-{Stealth[length=1.8mm]}] (O) -- (E3) node[below left=-1pt,text=FigureInk] {$e_3$};
  \fill[FigureStroke] ($1.0*(E1)$) circle (1.1pt);
  \fill[FigureStroke] ($0.5*(E2)$) circle (1.1pt) node[right,text=FigureInk] {$Ae_2$};
  \fill[FigureStroke] ($0.3333*(E3)$) circle (1.1pt) node[above left=-1pt,text=FigureInk] {$Ae_3$};
  \node[anchor=north] at (0.2,-0.95) {(a) 正交方向与它们的像};
\end{scope}
\begin{axis}[at={(34mm,-14mm)},anchor=south west,width=54mm,height=40mm,
  xmin=0.5,xmax=15.5,ymin=0,ymax=1.6,
  axis x line=bottom,axis y line=left,
  xtick={1,5,10,15},ytick={0,0.5,1,1.4142},yticklabels={$0$,$0.5$,$1$,$\sqrt2$},
  xlabel={$n$},ylabel={距离},
  xlabel style={font=\figurefont\scriptsize},ylabel style={font=\figurefont\scriptsize},
  tick label style={font=\figurefont\scriptsize,text=FigureInk},
  ymajorgrids,grid style={draw=FigureGrid,line width=.4pt},
  legend style={at={(0.97,0.62)},anchor=east,font=\figurefont\scriptsize,cells={anchor=west}},
  title={(b) 相邻输出之间的距离},title style={font=\figurefont\footnotesize}]
\addplot[FigureBlue,line width=1.1pt,mark=*,mark size=1pt] coordinates {(1,1.4142) (15,1.4142)};
\addplot[FigureStroke,line width=1.1pt,mark=*,mark size=1pt] coordinates {(1,1.1180) (2,0.6009) (3,0.4167) (4,0.3202) (5,0.2603) (6,0.2195) (7,0.1898) (8,0.1672) (9,0.1495) (10,0.1351) (11,0.1233) (12,0.1134) (13,0.1050) (14,0.0977) (15,0.0914)};
\legend{$\|e_{n+1}-e_n\|$,$\|Ae_{n+1}-Ae_n\|$}
\end{axis}
\end{tikzpicture}
